\section{Consequences and comparison with existing theory}\label{sec:scope}
The acquired law, its multiscale resolution and causal error propagation are distinct objects. Theorem~\ref{thm:bridge} connects them by domination of finite error measures. Proposition~\ref{prop:recharge-flow} derives this domination from raw expansion and refresh flows, without normalizing a stationary covariance by rare terminal probabilities. The countable singular realization verifies all these hypotheses, including a profile for which no limiting dimension exists. The finite anisotropic theorem, the stationary resolvent, the binary filter and the absorbing sensor retain mechanisms not subsumed by that particular refresh partition.

Two resource results address different questions. Theorem~\ref{thm:soft-state} quantifies retained information needed by an approximate continuation challenge and a geometric readout. Theorem~\ref{thm:single-probe} gives a one-task acquisition--resolution interaction. The independently selected-menu robust family remains a simultaneous lower witness with its own exact deficiency calculation, not a proof of an unrestricted law for every physical experiment.

\subsection{Theorem-level comparison}
\paragraph{Predictive states and sufficient statistics.}
Shalizi and Crutchfield's minimality and uniqueness theorems for causal states \cite[Theorems 2--3]{shalizi} identify minimal prescient representations of a stochastic process. Predictive state representations \cite{psr} use future action--observation tests. Our quotient proposition does not claim priority for either construction. It explicitly includes controlled continuation legality, paid resources and conditional-version domains. Its extra measurable work is the density-instrument descent at every declared continuous report. The categorical treatment of kernels and sufficiency in \cite{fritz} supplies another general language; an abstract sufficient statistic alone does not provide the acquired-mass or finite-machine estimates used here.

\paragraph{Bisimulation and belief-state control.}
Behavioral metrics for Markov decision processes \cite{ferns} compare states through rewards and transition laws and control value differences. Our metric is fixed by executable prediction tasks and valid update geometry; it is not asserted to be a universal bisimulation metric. Saldi, Y\"uksel and Linder \cite[Assumption 1 and Theorem 3]{saldi} obtain near-optimal finite belief-model policies for discounted partially observed control under their continuity, compactness and moment hypotheses. That theorem optimizes a controller, whereas our matching lower is for a fixed exploration and the law it actually acquires. Our finite-horizon two-sided resolution estimate does not imply their unrestricted control conclusion, nor does asymptotic policy near-optimality supply our actual-mass lower bound.

\paragraph{Quantization and nonlinear filtering.}
Conditional-mean projection and optimal probability quantization are classical \cite{graf}. Nonlinear filter approximation by quantization is developed in \cite{pages,sellami,feuer}. The additional obligation here is a matching lower for the very same acquired task, together with the forced-moment condition needed by a persistent finite machine. Zhu's local estimates for Ahlfors--David measures \cite{zhu}, and for Moran measures \cite{zhu-moran}, concern fine uniformity of contributions of optimal cells and successive quantization errors. The cylinder cover in Lemma~\ref{lem:moran-profile} is not offered as a replacement for that theory. Its role is to verify a multiscale profile, then propagate that profile through an actual nonstationary recurrent experiment with countably many weighted components. The finite-budget anisotropic results additionally retain width collapse and finite intersecting mixtures. We do not classify arbitrary singular measures.

\paragraph{Positive moment operators.}
Average contraction of iterated random functions \cite{diaconis} and moment stability of Markov jump systems \cite{costa} are established subjects. Ogura and Martin \cite{ogura} characterize mean stability for positive semi-Markovian jump systems with resets by a spectral-radius condition. That homogeneous stability question is not the joint domination of forced error measures by changing acquired measures required here. We claim no novelty for a Neumann resolvent or the criterion $\rho(K_2)<1$. The role of Theorem~\ref{thm:kernel} is more specific: reverse-edge operators propagate the conditional first and mixed second moments of the \emph{rounding forcing}, and the resulting matrix is normalized by the actual acquisition weights. Exact suspension cancels in the forced profiles, not in the physical budget. This distinction is why a recurrent expanding kernel can still yield a matching finite-memory risk curve. In the nonstationary partitioned class, Theorem~\ref{thm:bridge} instead propagates forward error measures and obtains the terminal mass factor directly; Proposition~\ref{prop:recharge-flow} identifies the necessary refresh-versus-expansion balance used by that sufficient criterion. We do not claim that this sufficient condition is a necessary spectral classification of all raw kernels.

\paragraph{Comparison of experiments.}
Blackwell comparison \cite{blackwell}, Le Cam's deficiency theory \cite{lecam}, and Torgersen's systematic account \cite{torgersen} provide the classical common-decision-loss principle. The triangle inequality and total-variation contraction in Theorem~\ref{thm:morphism} are instances of that principle. The added structure is the finite-resource causal implementation, with explicit lifted controller classes, stopping indices, clocks and numerical interfaces. A forward simulator yields an upper risk transfer, not a lower bound relative to a different oracle. The mandatory-state and robust converses are proved directly in Theorems~\ref{thm:tag} and \ref{thm:robust}, rather than attributed to a forward deficiency bound.

\subsection{Necessary distinctions illustrated by exact examples}
First, the valid envelope $[0,1]$ equipped with a point mass and the same envelope equipped with a uniform law have respectively zero one-point distortion and strictly positive finite-label distortion. Hence a reachability envelope is not an acquired law. More generally a component of probability $w$ contributes with weight $w$, even if its geometric diameter is large.

Second, an identity update has no need for repeated quantization. If the initial state has already been replaced by a representative, the same representative can be retained indefinitely. The bound $Nr$ obtained by blindly injecting a covering radius at every hold is not an inherent resource requirement. The acquisition-weighted certificate corrects precisely this failure of accounting.

Third, a statistical garbling can increase memory needed to match its \emph{own} oracle. Observing a fair hidden bit exactly produces only posterior probabilities zero and one. Adding independent Gaussian noise to this observation produces a continuously varying posterior probability. Two labels reproduce the first oracle exactly, but no finite number reproduces the latter continuously distributed oracle exactly under square score. This is compatible with data processing because the baselines differ.

Fourth, let $P_N=\delta_0$ and $Q_N=(1-e^{-Nc})\delta_0+e^{-Nc}\delta_1$, $c>0$. Their total variation distance tends to zero, but the exponential rates at the point one are infinite and $c$, respectively. The common-risk defect of this paper is therefore not a certificate of large-deviation equivalence.

Finally, a hidden mode or an internal simulator register is not free because its name is omitted from a state vector. Independent pairs of retained states have the product cardinality in Theorem~\ref{thm:morphism}. Conversely that product is a constructive upper bound, not a theorem that every simulator requires it. Theorem~\ref{thm:tag} establishes exact support necessity through zero-error continuation challenges; Theorem~\ref{thm:soft-state} replaces it by an information lower bound when errors are allowed. An exact label machine and a finite-bit numerical implementation are also different objects: the latter retains the nonzero calibration, workspace, and readout terms of \eqref{eq:joint}. Exact-arithmetic zero distortion is not a promise of a finite-precision physical output with zero error.


\subsection{Uniformity and the remaining mathematical questions}
The terminal-mass theorem permits nonstationarity and countably many scale-separated components. Its raw-flow realization permits arbitrarily rare components, changes of their acquired weights, recurrent expansion and exact waiting. This closes the rare-component normalization problem for the stated refresh-preserved partition class. It does not cover arbitrary transfers of a nonzero old approximation error between different strata; such transfers need a matrix-valued balance or another localization. Nor does the countable theorem handle arbitrary intersecting accumulations without its relative separation and small-ball certificates. The finite intersecting-stratum result remains independently valid.

The noisy continuation converse is not a matching theorem at every tag-error level. There is still no general shortest-program, minimum transient-workspace or optimal-runtime characterization below the constructive feasible regions. The one-probe theorem couples acquisition and resolution through a common scalar target, but does not establish an optimal joint exploration and calibration policy for a recurrent unknown-kernel model. Report thinning, complete causal deficiency and the exact terminal packet deficiency remain distinct quantities.

Learning unknown kernels with optimally charged pilots and unrestricted adaptive exploration remain substantive questions. Deterministic-horizon error-measure bounds are not automatically stopped bounds: a stopping rule may select atypically large error. The common bounded-stopping law comparisons already proved have their stated controller-uniform hypotheses. Infinite-horizon transcript comparison, large-deviation equivalence and unbounded-generator closure require different estimates and are not inferred from finite-time total variation. These are specific extensions of the same experiment--geometry--resource problem, not reasons to replace it by a model-specific surrogate.
